\begin{saresult}{A.120}{Proved under stated hypotheses}{Stable components without a smoothing hypothesis}
Let $M$ be any compact PL manifold, with the side boundary fixed in
the definition of concordance, and let $K$ be a finite simplicial
homotopy model for $M$. The complete countable presentation $E(K)$
of A.119 computes its stable PL concordance component group:
\[
 \operatorname*{colim}_{j}\pi_0P_{\rm PL}(M\times I^j)
       \cong E(K).
\]
There is no smoothability or dimension hypothesis on this stable
statement. It does not identify the unstable group with the stable
group for a nonsmoothable PL manifold.
\end{saresult}
\begin{proof}
In this proof Whitehead symbols denote the connective spectra,
rather than only their zeroth infinite-loop spaces. The definitions
in \cite[Definition~1.3.2 and (1.3.3), pp.~16--17;
Definition~1.3.4 and (1.3.5), p.~18]
{WaldhausenJahrenRognes2013} give
\[
\begin{split}
 \operatorname{Wh}^{\rm Diff}(K)
   &=\operatorname{cofib}(\Sigma^\infty K_+\longrightarrow A(K)),\\
 \operatorname{Wh}^{\rm PL}(K)
   &=\operatorname{cofib}(A(*)\wedge K_+\longrightarrow A(K)).
\end{split}
\]
The first arrow factors through the second by the unit
$\mathbb S\to A(*)$. Replace this composable pair by cofibrations
$U\hookrightarrow V\hookrightarrow W$ without changing its
homotopy cofibers. The quotient sequence
$V/U\to W/U\to W/V$ is a cofiber sequence. Therefore the
precise comparison sequence is
\[
 \operatorname{Wh}^{\rm Diff}(*)\wedge K_+
 \longrightarrow\operatorname{Wh}^{\rm Diff}(K)
 \longrightarrow\operatorname{Wh}^{\rm PL}(K).
                                                        \tag{A.120.1}
\]
This derives the comparison from the defining assembly maps and
does not presume that $M$ has a smoothing.

The spectrum $\operatorname{Wh}^{\rm Diff}(*)$ is $2$-connected.
It is connective, the unit is an isomorphism on rank $\pi_0$, and
$\pi_1\operatorname{Wh}^{\rm Diff}(*)=
\operatorname{Wh}_1(1)=0$; see
\cite[Theorem~0.3 and the following paragraph, p.~4]
{WaldhausenJahrenRognes2013}. To check degree $2$ without assuming
an extra algebraic $K$-group calculation, take a smooth disk $D^m$
with $m\geq9$. The group comparison of A.119 identifies
\[
 \pi_2\operatorname{Wh}^{\rm Diff}(*)
 \cong\operatorname*{colim}_{j}
           \pi_0P_{\rm Diff}(D^m\times I^j).
\]
After rounding corners each base is a disk. Burghelea--Lashof's
component comparison, already checked in A.103, identifies each
term with its PL term. Those PL disk concordance component groups
are zero by the Alexander trick, as explicitly used in
\cite[proof of Corollary~1.4.2, p.~23]
{WaldhausenJahrenRognes2013}. Thus the displayed group is zero.

The finite cell filtration of $K_+$ builds the smash product on the
left of (A.120.1) from nonnegative suspensions of a $2$-connected
spectrum. Its homotopy groups in degrees at most $2$ vanish.
The long exact sequence of (A.120.1) consequently gives a natural
group isomorphism
\[
 \pi_2\operatorname{Wh}^{\rm Diff}(K)
       \xrightarrow{\sim}\pi_2\operatorname{Wh}^{\rm PL}(K).
\]
The left side is $E(K)$ by A.119. On the right, the PL version of
the stable parametrized h-cobordism theorem and the bar-space
group-law argument in A.119 identify it with the stable PL
concordance component group. Both comparisons are isomorphisms of
groups. No unstable PL stabilization theorem has been used.
\end{proof}

\paragraph{Coefficient fields and the remaining unstable question.}
The same stable PL group is obtained from a realization over any
real closed field. One can use the relative version of A.115 at
each finite stage. Likewise the component argument in A.113
applies with the entire subpolyhedron
$M\times\{0\}\cup\partial M\times I$ fixed, so the stagewise PL
and definable component groups coincide. The standard PL positive
suspension can be chosen with fixed rational polyhedral bends;
these comparisons commute with it. This also identifies the
stable definable component group with $E(K)$. For a particular
unstable nonsmoothable manifold, the remaining precise issue is
whether its map to this stable group is bijective; A.120 supplies
no such assertion.
